% ==========================================================
% 8. DISEÑO METODOLÓGICO (Fase 5)
% 8.1 Diseño | 8.2 Operacionalización | 8.3 Protocolo de rigor
% 8.4 Población y muestra | 8.5 Técnicas, instrumentos y validación
% ==========================================================

\section{Diseño metodológico}

\subsection{Enfoque, alcance y diseño}

El estudio tiene un enfoque cuantitativo y un alcance correlacional (sección 4). Para
llevarlos a la práctica se adopta un diseño no experimental transversal: las variables se
observan tal como ocurren, en un solo momento y sin manipular ninguna, mediante una
encuesta aplicada una única vez a los estudiantes [11].

Se descartan los diseños experimental y cuasi-experimental porque el estudio no
manipula la exposición ni asigna estudiantes a grupos: pregunta si existe relación, no
cuánto cambia el efecto con una intervención. Se descarta el diseño longitudinal porque el
semestre no alcanza para observar a los mismos estudiantes en más de un momento. Se
descartan los diseños cualitativos porque las hipótesis (sección 5) anticipan relaciones
entre variables medibles.

\subsection{Operacionalización de variables}

La Tabla~\ref{tab:operacionalizacion} completa la definición nominal-conceptual de la
sección 6 con el indicador, el instrumento, la escala y la unidad de cada variable.

\bigskip
\begin{center}
\small
\setlength{\tabcolsep}{3pt}
\begin{longtable}{L{2.6cm} L{3.6cm} L{3.8cm} L{2cm} L{2.4cm}}
\caption{Matriz de operacionalización de variables}\label{tab:operacionalizacion} \\
\toprule
\textbf{Variable} & \textbf{Indicador} & \textbf{Instrumento} & \textbf{Escala} & \textbf{Unidad} \\
\midrule
\endfirsthead
\toprule
\textbf{Variable} & \textbf{Indicador} & \textbf{Instrumento} & \textbf{Escala} & \textbf{Unidad} \\
\midrule
\endhead
Tiempo de exposición a pantallas (independiente) & Horas diarias promedio frente a pantallas en las jornadas de estudio & Ítem de autorreporte en la encuesta, por tipo de dispositivo & De razón & Horas por día \\
\midrule
Hábitos posturales (independiente) & Puntaje del índice de hábitos posturales inadecuados & Sección de la encuesta con ocho ítems de frecuencia (ver numeral 8.5) & Ordinal (puntaje) & Puntos, de 0 a 32 \\
\midrule
Fatiga visual (dependiente) & Puntaje de severidad del síndrome visual informático & Cuestionario CVS-Q, con 16 síntomas evaluados por frecuencia e intensidad [21], [22] & De intervalo (puntaje) & Puntos, de 0 a 32 \\
\midrule
Molestias osteomusculares (dependiente) & Presencia de molestias por región corporal y número de regiones afectadas & Cuestionario Nórdico estandarizado [23], limitado a las regiones del estudio & Nominal por región; de razón para el conteo & Sí/no; regiones, de 0 a 6 \\
\midrule
Sexo (interviniente) & Sexo declarado & Ítem de la encuesta & Nominal & Femenino / masculino \\
\midrule
Semestre cursado (interviniente) & Semestre en que está matriculado & Ítem de la encuesta & Ordinal & Semestre (1 a 10) \\
\midrule
Dispositivo principal (interviniente) & Equipo que usa con mayor frecuencia para estudiar & Ítem de la encuesta & Nominal & Portátil / escritorio / otro \\
\bottomrule
\end{longtable}
\end{center}

\noindent El puntaje del CVS-Q se calcula multiplicando la frecuencia (0 a 2) por la
intensidad de cada síntoma y reclasificando el resultado en 0, 1 o 2 puntos; la suma de
los 16 síntomas va de 0 a 32, y un puntaje de 6 o más se interpreta como indicativo del
síndrome, como se aplica en [13]. Además del puntaje, se registra esa clasificación
dicotómica y la exposición de cuatro horas o más al día, que es el punto de corte usado en
[13], para comparar con la literatura. La escala de cada variable determina las pruebas
estadísticas de la Fase 6.

\subsection{Protocolo de rigor}

Por ser un estudio correlacional no hay grupo de tratamiento ni grupo de comparación: el
control de lo que podría confundir los resultados se hace de forma estadística. Las
variables intervinientes (sexo, semestre y dispositivo principal) se miden en todos los
participantes y se incluyen como covariables en el análisis de la Fase 6, en vez de
emparejar o bloquear estudiantes.

El protocolo es el mismo para todos los participantes: la misma encuesta digital, en el
mismo periodo de aplicación de dos semanas, con las mismas instrucciones y la misma
declaración de consentimiento. La Tabla~\ref{tab:amenazas} registra las amenazas a la
validez que se anticipan y cómo se tratan.

\bigskip
\begin{center}
\small
\begin{longtable}{L{2.4cm} L{5.4cm} L{6.8cm}}
\caption{Amenazas a la validez y su tratamiento}\label{tab:amenazas} \\
\toprule
\textbf{Tipo} & \textbf{Amenaza} & \textbf{Tratamiento} \\
\midrule
\endfirsthead
\toprule
\textbf{Tipo} & \textbf{Amenaza} & \textbf{Tratamiento} \\
\midrule
\endhead
Interna & Sesgo de información y de memoria, porque el tiempo de pantalla y los síntomas se autorreportan [13]. & Instrumentos validados para los síntomas, preguntas con periodos de referencia definidos y respuestas anónimas. Queda sin resolver: no hay medición objetiva. \\
\midrule
Interna & Variables que confunden la relación (sexo, semestre, dispositivo). & Se miden y se controlan estadísticamente. Quedan sin controlar los factores filtrados (iluminación, defectos refractivos, trauma previo). \\
\midrule
Interna & Imposibilidad de establecer causalidad, por el diseño transversal. & Las conclusiones se redactan como asociación y no como causa. \\
\midrule
Externa & Muestra por conveniencia, de un solo programa y una sola sede. & Se declara la limitación y no se generalizan los resultados a otros programas ni sedes. \\
\midrule
Externa & Sesgo de selección: pueden responder más quienes tienen molestias. & Se invita a todos los estudiantes sin mencionar los síntomas en la convocatoria; se reporta la tasa de respuesta. \\
\bottomrule
\end{longtable}
\end{center}

\subsection{Población y muestra}

La población son los estudiantes de pregrado matriculados en el programa de Ingeniería de
Sistemas y Computación de la UPTC, sede Tunja, durante el periodo de aplicación. La
unidad de análisis es cada estudiante, que responde una encuesta. Se incluyen los
estudiantes mayores de 18 años que aceptan el consentimiento informado y se excluyen las
encuestas incompletas.

\pendiente{Registrar el número de estudiantes matriculados en el programa en el periodo de aplicación, con el dato de la Escuela o de Admisiones.}

El muestreo es no probabilístico por cuotas según el semestre, con invitación abierta a
todos los estudiantes del programa. No es posible un muestreo aleatorio simple por no
contar con una lista de contacto de toda la población, lo que limita la generalización
(Tabla~\ref{tab:amenazas}).

El tamaño de la muestra se calculó para detectar una correlación de 0,20 entre dos
variables, con un nivel de significancia de 0,05 bilateral y una potencia de 0,80. Con la
transformación de Fisher, $n=\left[\frac{Z_{1-\alpha/2}+Z_{1-\beta}}{C}\right]^{2}+3$ con
$C=\frac{1}{2}\ln\frac{1+r}{1-r}=0{,}2027$, resulta $n=194$. Se prevé una pérdida del 10~\%
por encuestas incompletas, por lo que la meta es de 216 estudiantes. Si la población
resulta menor que esa cifra, la meta se ajusta al total de estudiantes disponibles y se
declara. Como referencia, el estudio [13] calculó una muestra de 204 personas para
comparar la exposición con el síndrome visual informático.

\subsection{Técnicas, instrumentos y validación}

La técnica es una encuesta autoadministrada en un formulario digital. Las variables de
percepción y autorreporte se capturan con cuestionarios, y no hay registros automáticos ni
observación directa. La Tabla~\ref{tab:instrumento} resume la estructura de la encuesta.

\bigskip
\begin{center}
\small
\begin{longtable}{c L{4.2cm} L{7.8cm}}
\caption{Estructura de la encuesta}\label{tab:instrumento} \\
\toprule
\textbf{N.°} & \textbf{Sección} & \textbf{Contenido} \\
\midrule
\endfirsthead
\toprule
\textbf{N.°} & \textbf{Sección} & \textbf{Contenido} \\
\midrule
\endhead
1 & Consentimiento informado & Finalidad académica, voluntariedad, anonimato y tratamiento de datos de salud \\
2 & Datos generales & Sexo, semestre y dispositivo principal de estudio \\
3 & Exposición a pantallas & Horas diarias frente a pantallas, por tipo de dispositivo \\
4 & Hábitos posturales & Ocho ítems propios, de frecuencia \\
5 & Fatiga visual & CVS-Q, 16 síntomas \\
6 & Molestias osteomusculares & Cuestionario Nórdico, seis regiones \\
\bottomrule
\end{longtable}
\end{center}

Para la fatiga visual se usa el CVS-Q, que evalúa la frecuencia y la intensidad de 16
síntomas y tiene propiedades psicométricas aceptables en su validación original [21]; se
usa su versión adaptada para Colombia, validada en una muestra de 300 trabajadores
usuarios de computador de una universidad pública, con consistencia interna de 0,88 [22].
Para las molestias osteomusculares se usa el Cuestionario Nórdico estandarizado [23], que
también se empleó en estudiantes universitarios [1].

Los dos instrumentos se validaron en trabajadores o en otras poblaciones, no en
estudiantes de ingeniería de sistemas, por lo que se reporta la consistencia interna del
CVS-Q en la muestra del estudio y se declara la limitación. Su uso requiere confirmar los
términos de uso de los autores.

\pendiente{Confirmar con los autores o con el sitio oficial los términos de uso del CVS-Q (Colombia), que es un instrumento con derechos reservados.}

El índice de hábitos posturales es un instrumento propio de ocho ítems, cada uno con
frecuencia de 0 (nunca) a 4 (siempre) y redactado para que un puntaje mayor indique hábitos
más inadecuados: (1) mantener la cabeza inclinada hacia la pantalla, (2) espalda sin
apoyo, (3) permanecer más de una hora sin cambiar de posición, (4) usar el portátil en la
cama o el sofá, (5) pantalla por debajo del nivel de los ojos, (6) no hacer pausas activas,
(7) no hacer estiramientos y (8) silla o mesa inadecuadas. Estos ítems se apoyan en los
factores asociados en [1] y [6].

Antes de aplicar la encuesta se realiza una prueba piloto con un grupo pequeño de
estudiantes ajenos a la muestra, para detectar preguntas ambiguas, el tiempo de respuesta y
problemas técnicos del formulario. La validez de contenido de los ítems propios se revisa
con dos docentes expertos, y la confiabilidad del índice se estima con el alfa de Cronbach
[24], con un valor de 0,70 o más como criterio de aceptación. La validez de constructo
se sustenta en la coherencia con el marco teórico (sección 7) y en el uso de instrumentos
ya validados para los otros dos grupos de síntomas.

\pendiente{Registrar el número de estudiantes de la prueba piloto, el tiempo de respuesta, los ajustes hechos a los ítems y el alfa de Cronbach obtenido.}
